\documentclass[11pt]{article}
\usepackage[margin=1in]{geometry}
\usepackage[T1]{fontenc}
\usepackage[utf8]{inputenc}
\usepackage{amsmath,amssymb,amsthm}
\usepackage{booktabs}
\usepackage{array}
\usepackage{enumitem}
\usepackage{xcolor}
\usepackage{hyperref}
\usepackage{url}
\hypersetup{colorlinks=true,linkcolor=blue!50!black,citecolor=blue!50!black,urlcolor=blue!50!black}
\setlist[itemize]{leftmargin=1.4em,itemsep=0.25em,topsep=0.25em}
\setlist[enumerate]{leftmargin=1.6em,itemsep=0.25em,topsep=0.25em}

\newtheorem{theorem}{Theorem}
\newtheorem{lemma}[theorem]{Lemma}
\newtheorem{proposition}[theorem]{Proposition}
\theoremstyle{definition}
\newtheorem{definition}[theorem]{Definition}
\newtheorem{axiom}[theorem]{Axiom}
\theoremstyle{remark}
\newtheorem{remark}[theorem]{Remark}

% Lean references are kept out of the main exposition; this footnote helper carries them.
\newcommand{\LeanFile}[1]{\footnote{\raggedright\textbf{Lean}: \path{#1}.\par}}
\newcommand{\R}{\ensuremath{\mathbb{R}}}
\newcommand{\N}{\ensuremath{\mathbb{N}}}
\newcommand{\Z}{\ensuremath{\mathbb{Z}}}

\sloppy

\title{Symmetry Foundations for Inference in Lean\\From Finite Lattices to Infinitary Probability}
\author{Zarathustra Amareis Goertzel\\\small Mettapedia Project}
\date{Extended abstract for JoeFest 2026}
\begin{document}
\maketitle

\begin{abstract}
We report a Lean~4 formalization of the algebraic route from order and symmetry to the
probability calculus and Shannon information measures, in the tradition of Knuth and
Skilling~\cite{KnuthSkilling2012}. The development makes precise which hypotheses each
representation theorem actually consumes: an Alimov-style \emph{no anomalous pairs}
condition (equivalent to a separation property) for the additive representation, a
positivity gate for the product representation, and an explicit regularity gate
(measurability or monotonicity) for the variational/Cauchy step. The same algebraic
skeleton extends to infinitary probability: under three independent completeness axioms
($\sigma$-closure of events, sequential completeness of the value scale, and Scott
continuity of the valuation), $\sigma$-additivity is recovered as a theorem. The route
thus offers a single auditable trail of hypotheses from ordered-semigroup symmetry,
through addition, multiplication, and entropy, to $\sigma$-additive probability.
\end{abstract}

\section{Introduction}

Foundations of probabilistic inference fall into two main camps. The
\emph{measure-theoretic} foundation of Kolmogorov~\cite{Kolmogorov1933} starts from a
$\sigma$-algebra of events and a countably additive probability measure as primitives.
The \emph{plausibility-theoretic} foundation of Cox~\cite{Cox1946} and
Jaynes~\cite{Jaynes2003} derives the sum and product rules from consistency desiderata on
a real-valued plausibility function, but is sensitive to richness assumptions on the
plausibility scale: Halpern~\cite{Halpern1999} showed that on small finite domains the
classical functional-equation arguments need not yield the probability calculus at all.

Knuth and Skilling~\cite{KnuthSkilling2012} occupy a third position. They start with
\emph{ordered} valuations on finite distributive lattices, identify a small list of
order-and-symmetry axioms, and derive addition, multiplication, conditioning, and entropy
as forced numerical representations. The route is finitary by design and avoids global
continuity hypotheses, so the algebraic argument never depends on a continuum of
plausibility values.

This extended abstract summarizes a Lean~4 formalization%
\LeanFile{Mettapedia/ProbabilityTheory/KnuthSkilling/FoundationsOfInference.lean}
of that route, and three further results obtained in the process: an Alimov-style necessity
witness for the density axiom, an explicit Cauchy regularity gate for the variational step,
and an infinitary bridge that recovers $\sigma$-additivity from the same algebraic
skeleton.

\section{Additive representation}

Let $S$ be a totally ordered set carrying a binary operation $\oplus$ that represents the
valuation of disjoint joins on a distributive lattice of events. Write $x^{(n)}$ for the
$n$-fold iteration $x \oplus \cdots \oplus x$.

\begin{axiom}[Ordered semigroup]
$\oplus$ is associative and strictly monotone in each argument: $x < y$ implies
$x \oplus z < y \oplus z$ and $z \oplus x < z \oplus y$.
\end{axiom}

\begin{axiom}[No anomalous pairs, Alimov~\cite{Alimov1950}]
There is no pair $a, b > 0$ with $a^{(n)} < b^{(n)} < a^{(n+1)}$ for every $n \ge 1$.
\end{axiom}

The negation describes \emph{infinitesimal} differences invisible to finite iteration;
ruling them out is precisely what forces one-dimensionality. An equivalent rational-density
formulation, which we call \emph{separation}, asks that for any $0 < x < y$ and any positive
$a$ there exist positive integers $m, n$ with $x^{(m)} < a^{(n)} \le y^{(m)}$. The
formalization proves the equivalence both in the bounded-below (probability) setting and in
a bilateral setting that distinguishes positive and negative cones, and shows that
separation already implies commutativity of $\oplus$%
\LeanFile{Mettapedia/ProbabilityTheory/KnuthSkilling/Core/Algebra.lean}.

\begin{theorem}[Additive representation]
\textbf{Assumptions.} $(S, \oplus)$ is a linearly ordered semigroup with strict
monotonicity and no anomalous pairs.
\textbf{Conclusion.} There exists an order embedding $\Theta : S \hookrightarrow (\R, +)$
with
\[
\Theta(x \oplus y) = \Theta(x) + \Theta(y),
\]
unique up to positive affine rescaling. If $S$ has an identity $0$ that is also the minimum
of the order, $\Theta$ can be normalized so that $\Theta(0) = 0$ and $\Theta(x) \ge 0$ for
all $x$.
\end{theorem}

We have three independent proof routes%
\LeanFile{Mettapedia/ProbabilityTheory/KnuthSkilling/Additive/}: H\"older embedding (the
conceptually cleanest path, via Eric Paul's \texttt{OrderedSemigroups}
library~\cite{Paul2024,Holder1901}), Dedekind cuts (an explicit density construction from
separation), and grid induction (a constructive finite-lattice route closer to the original
Knuth--Skilling exposition).

\begin{remark}[Necessity of the density axiom]
The lexicographic semigroup $(\N \times \N, +)$ satisfies associativity and strict
monotonicity but contains an anomalous pair: with $a = (1, 0)$ and $b = (1, 1)$, every
$a^{(n)} = (n, 0) < b^{(n)} = (n, n) < a^{(n+1)} = (n+1, 0)$. No order embedding into
$(\R, +)$ exists, so the density axiom is not a consequence of associativity and
monotonicity alone.
\end{remark}

\begin{remark}[Bounded-below normalization is also necessary]
$(\Z, +, \le)$ satisfies the ordered-semigroup axioms and has no anomalous pairs, yet every
order embedding $\Theta$ assigns negative values to negative integers%
\LeanFile{Mettapedia/ProbabilityTheory/KnuthSkilling/Additive/Counterexamples/NegativeWithoutIdentity.lean}.
Positivity of the representation is therefore not automatic from order preservation alone;
it requires that the identity element be the minimum of the order.
\end{remark}

\section{Product, conditioning, and the variational gate}

A second binary operation $\otimes$ represents the valuation of independent products. After
regrading $\oplus$ to addition, we ask that $\otimes$ distribute over $+$, be associative,
and be positive ($x, y > 0 \Rightarrow x \otimes y > 0$).

\begin{theorem}[Product representation]
Under the above assumptions on $\otimes$, there exists $C > 0$ such that
\[
x \otimes y = \frac{xy}{C}.
\]
\end{theorem}

The proof reduces to a Cauchy-type functional equation in one variable; positivity of
$\otimes$ supplies the monotonicity needed to exclude Hamel-basis solutions%
\LeanFile{Mettapedia/ProbabilityTheory/KnuthSkilling/Multiplicative/Main.lean}. From the
product representation and a chaining axiom on conditional valuations one then recovers
the familiar product and Bayes rules.

\paragraph{Variational potentials and entropy.}
For a variational potential $H$ whose derivative satisfies the separated functional
equation $H'(m_x m_y) = \lambda(m_x) + \mu(m_y)$, two distinct gates are
needed~\cite{Aczel1966,Kuczma2009}. \emph{Globality} requires that the local
Lagrange-multiplier identity hold across a sufficiently rich class of inputs (so that the
equation is global, not just stationary-point-local). \emph{Regularity} requires that the
resulting Cauchy equation be solved under measurability, monotonicity, or continuity at one
point, in order to exclude pathological Hamel solutions. With both gates in force,
$H'(m) = B + C \log m$, hence
\[
H(m) = A + Bm + C(m \log m - m),
\]
which specializes to Kullback--Leibler divergence and Shannon entropy:
\[
D(w \| u) = \sum_i w_i \log \frac{w_i}{u_i},
\qquad
S(p) = -\sum_i p_i \log p_i.
\]
Both gates are tracked explicitly in the formalization%
\LeanFile{Mettapedia/ProbabilityTheory/KnuthSkilling/Variational/Main.lean}; without
regularity, Cauchy's equation admits uncountably many non-measurable solutions and the
algebraic derivation of entropy fails. The same Cauchy step appears in the
Shore--Johnson axiomatization of cross-entropy~\cite{ShoreJohnson1980}, and it is at this
shared step that both routes must place their regularity gate.

\section[Infinitary bridge: sigma-additivity]{Infinitary bridge: $\sigma$-additivity}

A long-standing limitation of algebraic foundations of probability is that the resulting
valuation is only \emph{finitely} additive, while Kolmogorov's framework takes
$\sigma$-additivity as primitive. The de~Finetti--Kolmogorov debate~\cite{deFinetti1974}
is partly a debate about whether this gap is benign. The algebraic route closes the gap by
adding three natural extension axioms, each independent of the finitary symmetries.

\begin{axiom}[$\sigma$-complete events]
Every countable family $(e_n)$ of events in the lattice $E$ has a join $\bigsqcup_n e_n \in E$.
\end{axiom}

\begin{axiom}[Scale completeness]
Every bounded monotone sequence in the value scale $S$ has a supremum.
\end{axiom}

\begin{axiom}[Scott continuity of the valuation]
The valuation $v : E \to S$ preserves directed suprema: $v(\bigsqcup_i e_i) = \bigsqcup_i v(e_i)$.
\end{axiom}

\begin{theorem}[$\sigma$-additivity from K\&S\,+\,completeness]
Under the K\&S axioms (associativity, strict monotonicity, no anomalous pairs, bounded-below
identity) plus the three completeness axioms above, the composition $\mu = \Theta \circ v$
is $\sigma$-additive: for any pairwise disjoint sequence $(e_n) \subseteq E$,
\[
\mu\!\Bigl( \bigsqcup_n e_n \Bigr) \;=\; \sum_{n=0}^\infty \mu(e_n).
\]
\end{theorem}

The three axioms are mutually independent and none is a consequence of the finitary
axioms%
\LeanFile{Mettapedia/ProbabilityTheory/KnuthSkilling/Core/ScaleCompleteness.lean}: countable
joins can fail in a Boolean algebra that supports only finite ones; an Archimedean scale
need not be sequentially complete (the rationals); and a finitely-additive valuation need
not respect directed suprema (any purely finitely-additive measure on $\N$). Categorically,
the move from Boolean algebras to $\sigma$-frames is exactly the right setting for these
axioms, and Scott continuity is the natural morphism condition. The algebraic regrade
$\oplus \to +$ continues to do the structural work in the infinitary case:
$\sigma$-additivity becomes a theorem about the interaction of Scott continuity with the
additive coordinate, not a primitive postulate.

\section{Place in the landscape}

\begin{table}[h]
\centering
\footnotesize
\setlength{\tabcolsep}{4pt}
\begin{tabular}{lcccc}
\toprule
\textbf{Foundation} & \textbf{Events} & \textbf{Codomain} & \textbf{Key axiom} & \textbf{Regularity gate} \\
\midrule
Kolmogorov~\cite{Kolmogorov1933} & $\sigma$-algebra & $[0,1]$ & $\sigma$-additivity & --- \\
Cox--Jaynes~\cite{Cox1946,Jaynes2003} & propositions & $[0,1]$ & consistency & continuity \\
de~Finetti~\cite{deFinetti1974} & events & $[0,1]$ & finite additivity & --- \\
K\&S, finite~\cite{KnuthSkilling2012} & lattice & ordered scale & associativity & no anomalous pairs \\
K\&S\,+\,completeness (this work) & $\sigma$-lattice & complete scale & + Scott cont. & + completeness \\
Shore--Johnson~\cite{ShoreJohnson1980} & distributions & divergence & system independence & measurability \\
\bottomrule
\end{tabular}
\caption{Where each foundation places its primary structural axiom and its regularity gate.}
\end{table}

The Knuth--Skilling route makes finite additivity, multiplication, and entropy
\emph{derived} quantities; under the three completeness axioms above, $\sigma$-additivity
joins them as a theorem. Because the algebraic argument never invokes a continuum of
plausibility values, the finite-lattice case does not need to mimic one to support the
calculus, and each step from finite symmetry up to measure-theoretic probability appears
as a separately auditable bridge axiom rather than a primitive postulate.

\section*{Acknowledgments}
Thanks to Ben Goertzel for proposing the separation axiomatization and for first pointing
me toward the Knuth--Skilling paper, and to Eray \"Ozkural for ongoing discussion of the
algebraic foundations program. The formalization builds on Mathlib and on Eric Paul's
\texttt{OrderedSemigroups} library~\cite{Paul2024}.

\begin{thebibliography}{99}

\bibitem{KnuthSkilling2012}
K.~H. Knuth and J.~Skilling.
\newblock Foundations of inference.
\newblock \emph{Axioms}, 1(1):38--73, 2012.

\bibitem{Halpern1999}
J.~Y. Halpern.
\newblock Cox's theorem revisited.
\newblock \emph{Journal of Artificial Intelligence Research}, 11:429--435, 1999.

\bibitem{Cox1946}
R.~T. Cox.
\newblock Probability, frequency, and reasonable expectation.
\newblock \emph{American Journal of Physics}, 14:1--13, 1946.

\bibitem{Kolmogorov1933}
A.~N. Kolmogorov.
\newblock \emph{Foundations of the Theory of Probability}.
\newblock Chelsea, New York, 1956. Original German edition 1933.

\bibitem{Jaynes2003}
E.~T. Jaynes.
\newblock \emph{Probability Theory: The Logic of Science}.
\newblock Cambridge University Press, 2003.

\bibitem{deFinetti1974}
B.~de~Finetti.
\newblock \emph{Theory of Probability} (2 vols).
\newblock Wiley, 1974--1975.

\bibitem{Aczel1966}
J.~Acz\'el.
\newblock \emph{Lectures on Functional Equations and Their Applications}.
\newblock Academic Press, 1966.

\bibitem{Kuczma2009}
M.~Kuczma.
\newblock \emph{An Introduction to the Theory of Functional Equations and Inequalities:
Cauchy's Equation and Jensen's Inequality}, 2nd ed., edited by A.~Gil\'anyi.
\newblock Birkh\"auser, 2009.

\bibitem{Alimov1950}
N.~G. Alimov.
\newblock On ordered semigroups.
\newblock \emph{Izv. Akad. Nauk SSSR Ser. Mat.}, 14:569--576, 1950.

\bibitem{Holder1901}
O.~H\"older.
\newblock Die Axiome der Quantit\"at und die Lehre vom Mass.
\newblock \emph{Berichte \"uber die Verhandlungen der K\"oniglich S\"achsischen
Gesellschaft der Wissenschaften zu Leipzig, Math.-Phys. Cl.}, 53:1--64, 1901.

\bibitem{ShoreJohnson1980}
J.~E. Shore and R.~W. Johnson.
\newblock Axiomatic derivation of the principle of maximum entropy and the principle of
minimum cross-entropy.
\newblock \emph{IEEE Trans. Inform. Theory}, 26(1):26--37, 1980.

\bibitem{Paul2024}
E.~Paul.
\newblock OrderedSemigroups: Formalization in Lean~4.
\newblock GitHub repository, 2024. \url{https://github.com/ericluap/OrderedSemigroups}

\end{thebibliography}
\end{document}
